\SourceLeaf{72}%
einer gedrängten Übersicht über alle übrigen Schritte verbände, die zur ersten Bestimmung einer Kometenbahn nötig sind. Wir bemerken dazu, dass auch einige der anderen Formeln hier in einer Form erscheinen, die von der sonst üblichen abweicht. Diese Form bringt aber nur dann eine gewisse Abkürzung, wenn man die Tafel benutzt, die der Verfasser vor einem Jahr zur unmittelbaren Berechnung der Logarithmen von Summen und Differenzen veröffentlicht hat. So steht hier auf zwei oder drei Seiten alles beisammen, was man außer den Tafeln zur Berechnung einer Kometenbahn braucht. Zugleich wird alles an einem Beispiel zum letzten Kometen erläutert. Der Verfasser leitet aus seinen Beobachtungen vom 7., 14. und 21. April Elemente ab, die von den oben mitgeteilten und auf sämtliche Beobachtungen gestützten Elementen nur wenig abweichen.
\par\bigskip\centerline{\rule{.24\textwidth}{.3pt}}
